\documentclass[conference]{IEEEtran}
\usepackage[T1]{fontenc}
\usepackage{graphicx}
\usepackage{booktabs}
\usepackage{amsmath}
\usepackage{amssymb}
\usepackage{array}
\usepackage{url}
\usepackage{balance}
\usepackage[hidelinks]{hyperref}
\newcommand{\maybefigure}[3]{%
\begin{figure}[!t]\centering
\IfFileExists{#1}{\includegraphics[width=\columnwidth]{#1}}{%
\fbox{\parbox[c][1.05in][c]{0.92\columnwidth}{\centering Figure generated by \texttt{make reproduce}\\(Insert project figure when compiling locally)}}}%
\caption{#2}\label{#3}\end{figure}}
\title{From-Scratch K-Means Clustering of UNSW-NB15 Traffic:\\A Euclidean and Diagonal Mahalanobis Comparison}
\author{\IEEEauthorblockN{Jonathan Morse}\IEEEauthorblockA{CSEC 620 --- Cyber Analytics and Machine Learning\\Rochester Institute of Technology\\RIT ID: jam7142}}
\begin{document}
\maketitle
\begin{abstract}
This study implements Lloyd's K-means algorithm from scratch and compares Euclidean distance with a positive-diagonal Mahalanobis extension on the UNSW-NB15 network-traffic dataset. The reproducible experiment uses 3,730 class-balanced records and 39 standardized numeric features. Ground-truth traffic categories are excluded from centroid fitting and used for the prescribed sampling and evaluation. A fixed $k=11$, seed, K-means++ initialization, and 30-restart budget are shared across distance configurations. The predeclared primary criterion, silhouette evaluated in a common standardized Euclidean space, favors Euclidean K-means (0.4513 versus 0.4402). The weighted geometry nevertheless increases V-measure (0.2567 to 0.2758), majority-vote accuracy (0.3064 to 0.3303), and macro F1 (0.2347 to 0.2441), while reducing macro precision. We report both the tradeoffs and the possibility that exploratory label associations influenced feature-group selection. The results describe in-sample clustering, not independently validated intrusion detection.
\end{abstract}
\begin{IEEEkeywords}
K-means, UNSW-NB15, network intrusion, unsupervised clustering, Mahalanobis distance, silhouette
\end{IEEEkeywords}

\section{Introduction and Dataset}
Network-traffic clustering can reveal groups of flows with similar statistical and connection characteristics without requiring class labels during fitting. This project evaluates a student-implemented K-means algorithm on UNSW-NB15, a network-intrusion benchmark developed to represent normal traffic and contemporary attack behaviors \cite{moustafa2015}. The dataset distinguishes one normal category and nine attack categories. Its attributes span flow-level, basic, content, time, and connection-context information; this mix motivates examining whether feature weighting changes the discovered partitions.

The primary input is the course-provided \texttt{UNSW\_NB15\_training-set.csv} partition. The \texttt{attack\_cat} field is reserved for sampling and evaluation; \texttt{id} and the binary \texttt{label} field are removed. The supplied preprocessing pipeline retains 39 numeric features, excludes categorical fields, and standardizes the retained numeric columns. This makes features with different original units more comparable under squared-distance calculations.

The configured class-balanced sample cap is 4,000 records, but the actual sample has $n=3{,}730$: 400 records from each of nine categories and 130 Worms records (the available count). The ten categories include Normal and nine attacks. This sampling improves coverage of rare categories but changes class prevalence relative to the source data. Labels support only this prescribed sample construction and subsequent scoring; they are not included in $X$, centroid initialization, assignment, or centroid updates.

\section{From-Scratch K-Means}
\subsection{Initialization and Assignment}
The \texttt{KMeansScratch} implementation supports both random selection of $k$ distinct observations and K-means++ seeding. The reported runs use K-means++, in which the first centroid is sampled at random and later centroids are sampled with probabilities proportional to their squared distance from the nearest already-selected centroid \cite{arthur2007}. When all candidate distances vanish, the implementation chooses from remaining unselected observations.

Given observations $x_i\in\mathbb{R}^{39}$ and centroids $\mu_j$, the assignment step chooses
\begin{equation}
 z_i=\underset{1\leq j\leq k}{\arg\min}\;d^2(x_i,\mu_j).
\end{equation}
The squared-distance matrix is evaluated with vectorized NumPy operations; the cluster index is the minimum-distance column for each observation.

\subsection{Centroid Updates and Convergence}
For a nonempty cluster $S_j=\{i:z_i=j\}$, the update is
\begin{equation}
 \mu_j\leftarrow \frac{1}{|S_j|}\sum_{i\in S_j}x_i.
\end{equation}
If $S_j$ is empty, its previous centroid is retained instead of computing an undefined mean. Iterations terminate when either assignments do not change or the maximum Euclidean centroid displacement is below $10^{-4}$; the maximum is 300 iterations. The final assignments are recomputed after the last centroid update.

K-means can converge to different local optima. Accordingly, each primary run uses \texttt{n\_init=30} restarts with seed 42 and retains the restart with minimum within-geometry inertia,
\begin{equation}
 J=\sum_{i=1}^{n}d^2(x_i,\mu_{z_i}).
\end{equation}
The objective is used to select restarts \emph{within} a distance configuration, not to compare Euclidean and Mahalanobis scores across geometries.

\subsection{Implementation Validation}
The 42 supplied automated tests passed. The project grader also passed reproducibility, scratch-method completeness, absence of library K-means calls in the scratch class, and reference agreement in both geometries. Against a scikit-learn reference with matched geometry, Euclidean scratch/reference inertia ratio was 1.0139 (partition ARI 0.7550); Mahalanobis ratio was 1.0172 (ARI 0.7204). These are diagnostics, not the main clustering outcomes. Differences in random initialization and local optima can prevent identical partitions even when the objective and algorithm are correct.

\section{Choosing the Number of Clusters}
A sweep over $k=2,\ldots,15$ examined within-geometry inertia and silhouette, a measure of within-cluster cohesion relative to between-cluster separation \cite{rousseeuw1987}. The original Euclidean sweep with 10 restarts yielded its highest common-space silhouette at $k=11$ (0.4513). Its inertia fell from about 55,336 at $k=10$ to 43,754 at $k=11$, then only to about 43,125 at $k=12$, providing additional elbow-style evidence.

To investigate reference agreement and sensitivity to local minima, the restart budget was increased to 30. The final Euclidean sweep favored $k=10$ by silhouette; the final Mahalanobis sweep favored $k=13$. Nevertheless, the initially selected $k=11$ was retained as a fixed comparison condition rather than changing $k$ separately for each geometry. Thus $k=11$ is a defensible \emph{previously chosen comparison setting}, not the maximizer of the final 30-restart Euclidean sweep. The ten known categories do not prescribe $k$, since clustering is selected using internal geometric criteria rather than class count.

\maybefigure{../results/euclidean/k_selection.png}{Euclidean $k$-selection sweep (2--15) for the final configuration. The final sweep's silhouette optimum is $k=10$; the primary comparison retains the previously selected $k=11$.}{fig:k}

\section{Distance Extension and Experimental Design}
For Euclidean distance, $d_E^2(x,\mu)=\sum_f(x_f-\mu_f)^2$. The diagonal Mahalanobis extension uses
\begin{equation}
 d_C^2(x,\mu)=\sum_{f=1}^{39}c_f(x_f-\mu_f)^2,\qquad c_f>0.
\end{equation}
All weights are finite and positive. Equivalently, this is Euclidean squared distance after scaling feature $f$ by $\sqrt{c_f}$. Since the numeric inputs are standardized, the diagonal changes the \emph{relative contributions} of standardized feature differences rather than simply correcting raw units. Doubling a weight doubles that feature's squared-distance contribution for a given standardized difference; it does not double the distance itself.

The final diagonal sets weight 2 for seven connection-context/count attributes and weight 1 for the other 32 numeric attributes. The emphasized fields are \texttt{ct\_srv\_src}, \texttt{ct\_dst\_ltm}, \texttt{ct\_src\_dport\_ltm}, \texttt{ct\_dst\_sport\_ltm}, \texttt{ct\_dst\_src\_ltm}, \texttt{ct\_src\_ltm}, and \texttt{ct\_srv\_dst}. They summarize repeated source, destination, service, or port relationships, potentially capturing behavior not emphasized by uniform weighting. This is a hypothesis, not a guarantee of superior clusters.

\begin{table}[!t]
\caption{Final Diagonal in the 39-Feature Order}
\label{tab:weights}
\centering\footnotesize
\begin{tabular}{p{0.65\columnwidth}c}
\toprule
Feature indices in retained order & $c_f$\\
\midrule
28, 30, 31, 32, 33, 37, 38 & 2.0\\
All remaining indices (1--39) & 1.0\\
\bottomrule
\end{tabular}
\vspace{2pt}
\parbox{\columnwidth}{\footnotesize Indices 28, 30--33, 37, and 38 correspond respectively to \texttt{ct\_srv\_src}, \texttt{ct\_dst\_ltm}, \texttt{ct\_src\_dport\_ltm}, \texttt{ct\_dst\_sport\_ltm}, \texttt{ct\_dst\_src\_ltm}, \texttt{ct\_src\_ltm}, and \texttt{ct\_srv\_dst}. Full ordered names and vector appear in Appendix A.}
\end{table}

The primary criterion was declared before testing candidate weights: \emph{common-space Euclidean silhouette} (\texttt{silhouette\_euclidean}). Both partitions are scored in the same standardized Euclidean feature space, while \texttt{silhouette\_configured} reports silhouette under each run's own geometry. The two primary runs share the identical 3,730-point sample, $k=11$, seed 42, K-means++ seeding, standardization, 30 restarts, tolerance, and iteration cap. Only the distance geometry changes.

\section{Evaluation and Results}
The evaluation includes internal silhouette scores and external agreement measures using \texttt{attack\_cat}: homogeneity, completeness, V-measure, and majority-vote cluster-to-class accuracy, macro precision, macro recall, and macro F1. Majority-vote mapping assigns each discovered cluster its most frequent true class \emph{on the same sample}. Thus these scores describe in-sample alignment; they are not held-out detection accuracy. Confusion matrices visualize the cluster-to-class correspondence.

\begin{table}[!t]
\caption{Results at Fixed $k=11$ (Higher Is Better)}
\label{tab:results}\centering\footnotesize
\begin{tabular}{lrr}
\toprule
Metric & Euclidean & Mahalanobis\\
\midrule
Silhouette, common Euclidean & \textbf{0.4513} & 0.4402\\
Silhouette, configured geometry & \textbf{0.4513} & 0.4178\\
Homogeneity & 0.2070 & \textbf{0.2253}\\
Completeness & 0.3377 & \textbf{0.3557}\\
V-measure & 0.2567 & \textbf{0.2758}\\
Majority-vote accuracy & 0.3064 & \textbf{0.3303}\\
Macro precision & \textbf{0.3749} & 0.3169\\
Macro recall & 0.2915 & \textbf{0.3137}\\
Macro F1 & 0.2347 & \textbf{0.2441}\\
\bottomrule
\end{tabular}
\end{table}

\maybefigure{../results/euclidean/confusion_matrix.png}{Euclidean cluster-to-class confusion matrix at $k=11$. The plot reflects the same sample used to fit K-means and assign majority-vote labels.}{fig:euc}
\maybefigure{../results/mahalanobis/confusion_matrix.png}{Diagonal Mahalanobis cluster-to-class confusion matrix at $k=11$.}{fig:mah}

Table~\ref{tab:results} shows that Euclidean K-means wins under the predeclared criterion: common-space silhouette falls by 0.0111 under Mahalanobis weighting (0.4513 to 0.4402). The configured-geometry silhouette for Mahalanobis is lower still (0.4178), although that measure changes the evaluation geometry and is not used to select the primary winner. The weighted partition nevertheless improves homogeneity (0.2070 to 0.2253), completeness (0.3377 to 0.3557), V-measure (0.2567 to 0.2758), majority-vote accuracy (0.3064 to 0.3303), macro recall (0.2915 to 0.3137), and macro F1 (0.2347 to 0.2441). Macro precision declines from 0.3749 to 0.3169. The higher recall and lower precision indicate a tradeoff rather than uniform improvement in class agreement.

These outcomes are consistent with weighting connection-context variables changing the partition toward some known traffic distinctions while slightly reducing separation under the original equal-weight Euclidean metric. Because the dataset was standardized, the effect is best understood as reallocation of distance emphasis among features and their observed variation, not removal of an uncorrected scale disparity. Inertia values are deliberately \emph{not} compared between distance runs, since their objective functions have different weights.

\section{Exploratory Selection, Limitations, and Conclusion}
During preliminary exploration, associations between numeric features and known traffic categories were examined. Although labels were never supplied to the K-means fitting procedure, this exploration may have influenced the choice of which connection-context features received increased weight. The submitted diagonal is a fixed, positive feature-group specification rather than a formula computed from class labels, but the \emph{model-selection process} cannot be described as fully label-blind. Accordingly, external agreement improvements should be interpreted cautiously; the predeclared unsupervised common-space silhouette remains the primary criterion.

Further limitations include the deliberately balanced rather than naturally prevalent traffic sample; exclusion of categorical attributes; K-means sensitivity to initialization, $k$, and distance geometry; and the restriction to one simple positive diagonal. No independent validation partition was used to evaluate the majority-vote mapping or assess generalization. The study therefore makes no claim about deployed intrusion-detection performance. Future work could examine multiple seeds, a held-out evaluation design, alternative unsupervised feature-group hypotheses, and appropriate categorical-data treatment.

In conclusion, the from-scratch K-means implementation reproduces both distance configurations and passes the supplied validation checks. At fixed $k=11$, Euclidean clustering yields the stronger predeclared common-space silhouette, whereas the diagonal Mahalanobis variant modestly improves several in-sample class-agreement metrics but reduces macro precision. The experiment illustrates why internal cluster quality and external label agreement should be reported separately and why exploratory weight selection must be disclosed.

\appendices
\section{Ordered Numeric Features and Exact Diagonal}
The retained numeric features, in the order used by the diagonal, are:
\begin{enumerate}
\setlength{\itemsep}{0pt}\setlength{\parskip}{0pt}
\item \texttt{dur}\item \texttt{spkts}\item \texttt{dpkts}\item \texttt{sbytes}\item \texttt{dbytes}\item \texttt{rate}\item \texttt{sttl}\item \texttt{dttl}\item \texttt{sload}\item \texttt{dload}\item \texttt{sloss}\item \texttt{dloss}\item \texttt{sinpkt}\item \texttt{dinpkt}\item \texttt{sjit}\item \texttt{djit}\item \texttt{swin}\item \texttt{stcpb}\item \texttt{dtcpb}\item \texttt{dwin}\item \texttt{tcprtt}\item \texttt{synack}\item \texttt{ackdat}\item \texttt{smean}\item \texttt{dmean}\item \texttt{trans\_depth}\item \texttt{response\_body\_len}\item \texttt{ct\_srv\_src}\item \texttt{ct\_state\_ttl}\item \texttt{ct\_dst\_ltm}\item \texttt{ct\_src\_dport\_ltm}\item \texttt{ct\_dst\_sport\_ltm}\item \texttt{ct\_dst\_src\_ltm}\item \texttt{is\_ftp\_login}\item \texttt{ct\_ftp\_cmd}\item \texttt{ct\_flw\_http\_mthd}\item \texttt{ct\_src\_ltm}\item \texttt{ct\_srv\_dst}\item \texttt{is\_sm\_ips\_ports}
\end{enumerate}
The exact vector is
\begin{equation*}
\begin{split}
C=\operatorname{diag}(&1,1,1,1,1,1,1,1,1,1,\\
&1,1,1,1,1,1,1,1,1,1,\\
&1,1,1,1,1,1,1,2,1,2,\\
&2,2,2,1,1,1,2,2,1).
\end{split}
\end{equation*}

\begin{thebibliography}{3}
\bibitem{moustafa2015} N. Moustafa and J. Slay, ``UNSW-NB15: A comprehensive data set for network intrusion detection systems (UNSW-NB15 network data set),'' in \emph{2015 Military Communications and Information Systems Conference (MilCIS)}, 2015, pp. 1--6, doi: 10.1109/MilCIS.2015.7348942.
\bibitem{arthur2007} D. Arthur and S. Vassilvitskii, ``k-means++: The advantages of careful seeding,'' in \emph{Proc. 18th Annual ACM-SIAM Symposium on Discrete Algorithms (SODA)}, 2007, pp. 1027--1035.
\bibitem{rousseeuw1987} P. J. Rousseeuw, ``Silhouettes: A graphical aid to the interpretation and validation of cluster analysis,'' \emph{Journal of Computational and Applied Mathematics}, vol. 20, pp. 53--65, 1987, doi: 10.1016/0377-0427(87)90125-7.
\end{thebibliography}
\end{document}
